\documentclass[12pt,a4paper]{article}

\usepackage[utf8]{inputenc}
\usepackage[T1]{fontenc}
\usepackage[brazil]{babel}
\usepackage{amsmath}
\usepackage{booktabs}
\usepackage{geometry}
\usepackage{array}
\usepackage{float}

\geometry{
    left=3cm,
    right=2cm,
    top=3cm,
    bottom=2cm
}

\renewcommand{\thesection}{\arabic{section}.}
\renewcommand{\thesubsection}{\arabic{section}.\arabic{subsection}.}

\title{\textbf{Análise das Medidas de Tendência Central\\
Dados de Aluguel de Bicicletas}}
\author{}
\date{}

\begin{document}

\maketitle

\section{Introdução}

Este relatório apresenta uma análise das medidas de tendência central das
principais variáveis numéricas relacionadas ao aluguel de bicicletas em uma
cidade dos Estados Unidos.

As medidas analisadas são a média, a mediana e, quando aplicável, a moda.
Essas medidas permitem resumir o comportamento dos dados por meio de valores
representativos e ajudam a compreender características importantes da
utilização do sistema de compartilhamento de bicicletas.

O conjunto de dados analisado possui 300 observações.

\section{Seleção das variáveis numéricas relevantes}

Para a análise das medidas de tendência central, foram selecionadas as
variáveis numéricas que possuem significado quantitativo no contexto do
conjunto de dados.

As variáveis consideradas relevantes são:

\begin{itemize}

    \item \textbf{temp}: representa a temperatura observada no período
    correspondente a cada registro;

    \item \textbf{casual}: representa a quantidade de usuários casuais,
    isto é, usuários que utilizaram o serviço sem possuir registro
    permanente no sistema;

    \item \textbf{registered}: representa a quantidade de usuários
    registrados que utilizaram o sistema;

    \item \textbf{total\_user}: representa o número total de usuários do
    serviço, obtido pela soma dos usuários casuais e registrados.

\end{itemize}

A variável \texttt{total\_user} é definida por:

\begin{equation}
\text{total\_user}_i =
\text{casual}_i + \text{registered}_i
\end{equation}

em que $i$ representa cada observação do conjunto de dados.

Essas variáveis foram escolhidas porque possuem interpretação quantitativa.
A temperatura representa uma característica numérica das condições
ambientais, enquanto as outras três variáveis representam diretamente a
utilização do sistema de bicicletas.

A variável \texttt{instant} não foi considerada nesta análise porque funciona
como um identificador sequencial das observações. Apesar de ser armazenada
como número, seus valores não representam uma quantidade de interesse.

As variáveis \texttt{season} e \texttt{weathersit} também podem aparecer
originalmente representadas por números, porém esses números correspondem a
categorias. Dessa forma, calcular média ou mediana desses códigos não possui
interpretação adequada.

A variável \texttt{dteday}, por sua vez, representa uma data e, portanto,
também não é utilizada diretamente nos cálculos de tendência central
realizados neste relatório.

\section{Medidas de tendência central}

As medidas de tendência central utilizadas são média, mediana e moda.

\subsection{Média aritmética}

A média aritmética é obtida pela soma de todos os valores de uma variável,
dividida pelo número total de observações.

Sua fórmula geral é:

\begin{equation}
\bar{x} =
\frac{\sum_{i=1}^{n} x_i}{n}
\end{equation}

em que:

\begin{itemize}
    \item $x_i$ representa cada observação;
    \item $n$ representa o número total de observações;
    \item $\bar{x}$ representa a média.
\end{itemize}

Neste conjunto de dados:

\begin{equation}
n = 300
\end{equation}

\subsubsection*{Média da temperatura}

A soma de todos os valores da variável \texttt{temp} é:

\begin{equation}
\sum temp_i = 6344{,}1
\end{equation}

Aplicando a fórmula da média:

\begin{equation}
\overline{temp} =
\frac{6344{,}1}{300}
\end{equation}

\begin{equation}
\overline{temp} = 21{,}147
\end{equation}

Portanto, a temperatura média registrada no conjunto de dados é de
aproximadamente:

\begin{equation}
\boxed{21{,}15}
\end{equation}

\subsubsection*{Média dos usuários casuais}

A soma de todos os valores de usuários casuais é:

\begin{equation}
\sum casual_i = 223130
\end{equation}

Logo:

\begin{equation}
\overline{casual} =
\frac{223130}{300}
\end{equation}

\begin{equation}
\overline{casual} =
743{,}7667
\end{equation}

Portanto:

\begin{equation}
\boxed{\overline{casual} \approx 743{,}77}
\end{equation}

Isso significa que foram observados, em média, aproximadamente 744 usuários
casuais por registro.

\subsubsection*{Média dos usuários registrados}

A soma dos valores da variável \texttt{registered} é:

\begin{equation}
\sum registered_i = 830834
\end{equation}

Assim:

\begin{equation}
\overline{registered} =
\frac{830834}{300}
\end{equation}

\begin{equation}
\overline{registered} =
2769{,}4467
\end{equation}

Portanto:

\begin{equation}
\boxed{\overline{registered} \approx 2769{,}45}
\end{equation}

Isso significa que o número médio de usuários registrados foi de
aproximadamente 2769 usuários por observação.

\subsubsection*{Média do total de usuários}

A soma de todos os valores da variável \texttt{total\_user} é:

\begin{equation}
\sum total\_user_i = 1053964
\end{equation}

Aplicando a fórmula:

\begin{equation}
\overline{total\_user} =
\frac{1053964}{300}
\end{equation}

\begin{equation}
\overline{total\_user} =
3513{,}2133
\end{equation}

Portanto:

\begin{equation}
\boxed{\overline{total\_user} \approx 3513{,}21}
\end{equation}

Em média, o sistema apresentou aproximadamente 3513 usuários por
observação.

Os resultados das médias são resumidos na Tabela a seguir.

\begin{table}[H]
\centering
\caption{Média das variáveis numéricas analisadas}
\begin{tabular}{lc}
\toprule
\textbf{Variável} & \textbf{Média} \\
\midrule
temp & 21,147 \\
casual & 743,77 \\
registered & 2769,45 \\
total\_user & 3513,21 \\
\bottomrule
\end{tabular}
\end{table}

\subsection{Mediana}

A mediana corresponde ao valor central de um conjunto de dados após os
valores serem colocados em ordem crescente.

Como o conjunto possui 300 observações, o número de observações é par.

Nesse caso, a mediana é calculada usando os valores das posições:

\begin{equation}
\frac{n}{2}
\end{equation}

e

\begin{equation}
\frac{n}{2} + 1
\end{equation}

Como:

\begin{equation}
n = 300
\end{equation}

temos:

\begin{equation}
\frac{300}{2} = 150
\end{equation}

e:

\begin{equation}
150 + 1 = 151
\end{equation}

Assim, em todas as variáveis analisadas, a mediana é calculada pela média
entre o 150º e o 151º valores da sequência ordenada.

A fórmula utilizada é:

\begin{equation}
Md =
\frac{x_{150} + x_{151}}{2}
\end{equation}

\subsubsection*{Mediana da temperatura}

Após ordenar os valores de temperatura:

\begin{equation}
x_{150} = 22{,}2
\end{equation}

e:

\begin{equation}
x_{151} = 22{,}2
\end{equation}

Portanto:

\begin{equation}
Md_{temp} =
\frac{22{,}2 + 22{,}2}{2}
\end{equation}

\begin{equation}
Md_{temp} =
\frac{44{,}4}{2}
\end{equation}

\begin{equation}
\boxed{Md_{temp} = 22{,}2}
\end{equation}

\subsubsection*{Mediana dos usuários casuais}

Os valores localizados nas posições 150 e 151 são:

\begin{equation}
x_{150} = 673
\end{equation}

\begin{equation}
x_{151} = 673
\end{equation}

Assim:

\begin{equation}
Md_{casual} =
\frac{673 + 673}{2}
\end{equation}

\begin{equation}
Md_{casual} =
\frac{1346}{2}
\end{equation}

\begin{equation}
\boxed{Md_{casual} = 673}
\end{equation}

\subsubsection*{Mediana dos usuários registrados}

Após ordenar os valores de \texttt{registered}, temos:

\begin{equation}
x_{150} = 3084
\end{equation}

e:

\begin{equation}
x_{151} = 3102
\end{equation}

Aplicando a fórmula:

\begin{equation}
Md_{registered} =
\frac{3084 + 3102}{2}
\end{equation}

\begin{equation}
Md_{registered} =
\frac{6186}{2}
\end{equation}

\begin{equation}
\boxed{Md_{registered} = 3093}
\end{equation}

\subsubsection*{Mediana do total de usuários}

Para a variável \texttt{total\_user}, os dois valores centrais são:

\begin{equation}
x_{150} = 3982
\end{equation}

e:

\begin{equation}
x_{151} = 4010
\end{equation}

Logo:

\begin{equation}
Md_{total\_user} =
\frac{3982 + 4010}{2}
\end{equation}

\begin{equation}
Md_{total\_user} =
\frac{7992}{2}
\end{equation}

\begin{equation}
\boxed{Md_{total\_user} = 3996}
\end{equation}

Os valores das medianas são apresentados a seguir.

\begin{table}[H]
\centering
\caption{Medianas das variáveis analisadas}
\begin{tabular}{lc}
\toprule
\textbf{Variável} & \textbf{Mediana} \\
\midrule
temp & 22,2 \\
casual & 673 \\
registered & 3093 \\
total\_user & 3996 \\
\bottomrule
\end{tabular}
\end{table}

\subsection{Moda}

A moda corresponde ao valor que aparece com maior frequência em um conjunto
de dados.

Matematicamente, pode ser representada por:

\begin{equation}
Mo =
\underset{x}{\operatorname{arg\,max}}\; f(x)
\end{equation}

em que $f(x)$ representa a frequência de ocorrência do valor $x$.

Diferentemente da média e da mediana, uma variável pode possuir uma única
moda ou várias modas.

\subsubsection*{Moda da temperatura}

Na variável \texttt{temp}, o valor:

\begin{equation}
26{,}0
\end{equation}

aparece 6 vezes no conjunto de dados.

Nenhum outro valor possui frequência superior a 6.

Portanto:

\begin{equation}
\boxed{Mo_{temp} = 26{,}0}
\end{equation}

\subsubsection*{Moda dos usuários casuais}

Na variável \texttt{casual}, o valor:

\begin{equation}
639
\end{equation}

aparece 3 vezes.

Essa é a maior frequência encontrada na variável.

Portanto:

\begin{equation}
\boxed{Mo_{casual} = 639}
\end{equation}

\subsubsection*{Moda dos usuários registrados}

Na variável \texttt{registered}, a maior frequência observada é igual a 2.

Entretanto, diversos valores aparecem exatamente duas vezes:

\begin{equation}
\begin{split}
Mo_{registered} = \{&
674,\ 1368,\ 1506,\ 1628,\ 1707,\ 1730,\ 2419,\ 2549,\\
&3107,\ 3348,\ 3413,\ 3840,\ 3848,\ 3854,\ 3896,\ 4044
\}
\end{split}
\end{equation}

Todos esses valores possuem a mesma frequência máxima de 2 ocorrências.

Dessa forma, a distribuição é multimodal e não possui uma única moda.

\subsubsection*{Moda do total de usuários}

Para a variável \texttt{total\_user}, a frequência máxima encontrada também
é igual a 2.

Os valores com essa frequência são:

\begin{equation}
\begin{split}
Mo_{total\_user} = \{&
1685,\ 2077,\ 3351,\ 4274,\ 4401,\\
&4758,\ 5119,\ 5202,\ 5312
\}
\end{split}
\end{equation}

Portanto, a variável \texttt{total\_user} também possui comportamento
multimodal.

Os resultados das modas são resumidos na tabela a seguir.

\begin{table}[H]
\centering
\caption{Modas das variáveis analisadas}
\begin{tabular}{p{3cm}p{10cm}}
\toprule
\textbf{Variável} & \textbf{Moda} \\
\midrule
temp & 26,0, com frequência 6 \\
casual & 639, com frequência 3 \\
registered & 674, 1368, 1506, 1628, 1707, 1730, 2419, 2549,
3107, 3348, 3413, 3840, 3848, 3854, 3896 e 4044, todos com frequência 2 \\
total\_user & 1685, 2077, 3351, 4274, 4401, 4758, 5119, 5202 e 5312,
todos com frequência 2 \\
\bottomrule
\end{tabular}
\end{table}

\section{Resumo das medidas de tendência central}

Os resultados gerais são apresentados na tabela abaixo.

\begin{table}[H]
\centering
\caption{Resumo das medidas de tendência central}
\begin{tabular}{lccc}
\toprule
\textbf{Variável} &
\textbf{Média} &
\textbf{Mediana} &
\textbf{Moda} \\
\midrule
temp & 21,15 & 22,2 & 26,0 \\
casual & 743,77 & 673 & 639 \\
registered & 2769,45 & 3093 & Multimodal \\
total\_user & 3513,21 & 3996 & Multimodal \\
\bottomrule
\end{tabular}
\end{table}

\section{Interpretação dos resultados}

\subsection{Temperatura}

A temperatura média observada foi de aproximadamente 21,15, enquanto a
mediana foi de 22,2.

Isso significa que o centro da distribuição das temperaturas está próximo
de 22. A mediana indica que metade das observações apresenta temperatura
inferior ou igual a aproximadamente 22,2 e a outra metade apresenta valores
superiores ou iguais a esse valor.

A moda foi 26,0, que apareceu 6 vezes, sendo a temperatura individual mais
frequente no conjunto de dados.

Como a média é ligeiramente inferior à mediana, existem observações de
temperaturas mais baixas que contribuem para reduzir o valor médio.

\subsection{Usuários casuais}

A média de usuários casuais foi de aproximadamente 743,77 usuários, enquanto
a mediana foi de 673 usuários.

A média maior que a mediana indica que existem observações com números
relativamente elevados de usuários casuais que aumentam o valor médio.

A moda foi 639 usuários, com três ocorrências.

Isso mostra que, apesar de a média ser aproximadamente 744 usuários, um dos
valores mais representativos da região central do conjunto encontra-se em
torno de 639 a 673 usuários.

\subsection{Usuários registrados}

A média de usuários registrados foi de aproximadamente 2769,45, enquanto a
mediana foi de 3093.

Neste caso, a mediana é maior que a média. Isso indica que existem
observações com números relativamente baixos de usuários registrados que
reduzem o valor da média.

A diferença também mostra que a distribuição dos usuários registrados não é
perfeitamente simétrica.

A variável não possui uma única moda. Existem 16 valores diferentes que
aparecem duas vezes cada, que é a maior frequência observada.

Portanto, a moda possui pouca capacidade de representar um único valor
típico para essa variável.

\subsection{Total de usuários}

A média do total de usuários foi de aproximadamente 3513,21 usuários.

A mediana foi igual a 3996.

Isso significa que, após ordenar os 300 valores, o centro do conjunto está
próximo de 3996 usuários.

A mediana maior que a média indica que algumas observações com quantidades
relativamente baixas de usuários reduzem o valor médio do conjunto.

No contexto do sistema de aluguel de bicicletas, isso sugere que existem
períodos com baixa utilização do serviço, enquanto uma parcela significativa
das observações apresenta utilização superior à média geral.

A variável \texttt{total\_user} também é multimodal. Nove valores diferentes
aparecem duas vezes cada, e nenhuma quantidade aparece mais frequentemente
do que isso.

Por esse motivo, neste caso, a média e principalmente a mediana são mais
informativas para representar o comportamento central da utilização total
do serviço do que uma única moda.

\subsection{Comparação entre usuários casuais e registrados}

Os resultados mostram uma diferença importante entre os dois tipos de
usuários.

A média de usuários casuais foi:

\begin{equation}
743{,}77
\end{equation}

enquanto a média de usuários registrados foi:

\begin{equation}
2769{,}45
\end{equation}

A diferença entre essas médias é:

\begin{equation}
2769{,}45 - 743{,}77 = 2025{,}68
\end{equation}

Portanto, em média, existem aproximadamente 2026 usuários registrados a mais
que usuários casuais por observação.

Também é possível comparar a participação média de cada grupo em relação ao
total médio de usuários.

Para os usuários casuais:

\begin{equation}
\frac{743{,}77}{3513{,}21} \times 100
\approx 21{,}17\%
\end{equation}

Para os usuários registrados:

\begin{equation}
\frac{2769{,}45}{3513{,}21} \times 100
\approx 78{,}83\%
\end{equation}

Dessa forma, considerando os valores médios, aproximadamente 78,83\% dos
usuários são registrados e 21,17\% são usuários casuais.

Isso mostra que a utilização do sistema de aluguel de bicicletas é composta
majoritariamente por usuários registrados no período representado pelos
dados analisados.

\section{Considerações finais}

A análise das medidas de tendência central permitiu resumir o comportamento
das principais variáveis numéricas do conjunto de dados.

A temperatura apresentou média de 21,15, mediana de 22,2 e moda de 26,0.

A quantidade de usuários casuais apresentou média de 743,77, mediana de 673
e moda de 639.

Para os usuários registrados, a média foi de 2769,45 e a mediana de 3093.
A variável apresentou várias modas com frequência igual a 2.

O total de usuários apresentou média de 3513,21 e mediana de 3996, também
apresentando diversas modas.

Os resultados demonstram que os usuários registrados representam a maior
parcela da utilização do serviço, correspondendo a aproximadamente 78,83\%
do número médio total de usuários.

Além disso, as diferenças observadas entre média e mediana indicam que as
distribuições não são perfeitamente simétricas. Em especial, a mediana do
total de usuários é superior à média, indicando a presença de observações
com baixa utilização que reduzem o valor médio.

Essas medidas fornecem uma visão inicial do comportamento do sistema de
aluguel de bicicletas e podem servir de base para análises posteriores que
relacionem a demanda a fatores como temperatura, estação do ano, condições
meteorológicas e período de observação.

\end{document}